% chapters/introduction.tex

\chapter{Introduction}
\label{ch:introduction}

% ── Motivation ────────────────────────────────────────────────────
\section{Motivation}
\label{sec:intro-motivation}

Denotational semantics is an approach to formalizing the behaviour and 
structure of programs by assigning mathematical meaning to the syntax of 
programming languages. This usually takes the form of interpreting 
types as domains and terms as continuous functions between them.  
The mathematical backbone of this enterprise is
\emph{domain theory}---the study of $\omega$-complete partial orders
($\omega$-CPOs), Scott-continuous maps, and their fixed points---which
provides the structures needed to give meaning to recursion and
partiality \cite{AJ1994}.

A landmark in the mechanisation of these ideas is the \Rocq{}
(formerly Coq) library of Benton, Kennedy, and Varming
\cite{BKV2009}, which formalised $\omega$-CPOs, the Kleene
fixed-point theorem, and the denotational semantics of PCF in Coq~8.2.
While influential, that library predates modern proof-assistant
infrastructure: it uses hand-rolled typeclasses, lacks any notion of
enriched categories, and has not been maintained for over a decade.

Meanwhile with the development of quantum architecture, interest in 
\emph{quantum programming languages} has motivated a \emph{quantum} 
generalisation of domain theory.  Kornell, Lindenhovius, and Mislove 
\cite{KLM2024} replace the classical Boolean-valued order $x \sqsubseteq y$ with a
quantale-valued order $R(x,y) \in Q$, yielding \emph{quantum CPOs}
whose hom-sets are themselves CPOs---a form of CPO-enrichment that
arises naturally from the quantum setting.  Formalising this richer
structure demands an enriched-category layer that \cite{BKV2009} never
provided.

This thesis addresses both gaps.  We present a from-scratch
domain-theory library for \Rocq{} 9.0, built on the Hierarchy Builder
framework \cite{HB2023}, that
%
\begin{enumerate}
  \item modernises and extends the classical development of
        \cite{BKV2009} (partial orders through the Kleene theorem,
        products, sums, function spaces, the lift monad, and a full
        PCF adequacy proof);
  \item introduces a CPO-enriched category formalisation (enriched
        categories, locally continuous functors, natural
        transformations, and a term-level Yoneda lemma); and
  \item lays the groundwork for quantum domain theory by formalising
        involutive quantales, quantum posets, quantum CPOs, and
        the quantum Kleene fixed-point theorem following \cite{KLM2024}.
\end{enumerate}
%
The enriched-category layer is not merely an afterthought: it is the
organising principle that connects the classical and quantum theories,
since both $\CPO$ and $\qCPO$ are instances of CPO-enriched
categories.

% NOTE: You can cross-reference later chapters like this:
%   \cref{ch:background}  →  "Chapter 2"
%   \cref{sec:bg-cpo}     →  "Section 2.1.2"
% cleveref handles the "Chapter"/"Section" prefix automatically.

% ── Contributions ─────────────────────────────────────────────────
\section{Contributions}
\label{sec:intro-contributions}

% TODO: List the main contributions:
% 1. Modernized domain theory library (~15,000 lines, Rocq 9.0 + HB)
% 2. CPO-enriched category formalization (enriched cats, LC functors, nat trans, Yoneda)
% 3. Recursive domain equations via EP-pairs and bilimits
% 4. Full PCF adequacy proof (soundness + adequacy via logical relation)
% 5. Discussion of quantum extensions (qCPO framework)

% ── Thesis outline ────────────────────────────────────────────────
\section{Thesis Outline}
\label{sec:intro-outline}

% TODO: Brief chapter-by-chapter summary
